\documentclass[11pt]{article}
\usepackage{fullpage,url,hyperref}
\usepackage{amsmath}
\usepackage{amssymb}

\usepackage[letterpaper,top=1in,bottom=1in,left=1in,right=1in,nohead]{geometry}

\setlength{\parindent}{0in}
\setlength{\parskip}{6pt}

\begin{document}
\thispagestyle{empty}
{\large{\bf ECE 6701 + CS 6501: Optimization for Engineering \hfill Due Wed, Oct
28}}\\

{\LARGE{\bf Homework 3: Duality}}
\vspace{0.2\baselineskip}
\hrule

{\bf Instructions:} Submit a single Jupyter notebook (.ipynb) of your work to
Canvas by 11:59pm on the due date. All math must be formatted using LaTeX. {\bf
  Be sure to show all the work involved in deriving your answers! If you just
  give a final answer without explanation, you may not receive credit for that
  question.}

You may discuss the concepts with your classmates, look on the internet for
help, or ask an LLM for advice, but first try hard to solve problems on your
own, and always write up the answers entirely on your own. Make sure you
understand the concepts and solutions---don't just look up the answers or use AI
to answer it for you. {\bf Cite any sources you used outside of the class
  material (webpages, LLMs, etc.).}

\begin{enumerate}

\item (Extended version of Exercise 5.3 from B\&V) Consider the optimization
  problem
  \begin{align*}
    \text{minimize} \quad & c^T x \\
    \text{subject to} \quad & f(x) \leq 0,
  \end{align*}
  with $c \neq 0$. We do not assume $f$ is convex.

  \begin{enumerate}
    \item Write down the dual problem in terms of the conjugate $f^*$. Explain
      why the dual problem is convex. (Don't just say it's because dual problems
      are always convex! Verify that the conditions for convexity are satisfied
      for this particular dual problem.)

    \item Write down the dual problem for the specific constraint:
      $$f(x) = x^T Q x - 1,$$
      where $Q \in S_{++}^n$.

    \item Give analytic solutions (in terms of $Q$ and $c$) for $x^*$ and
      $\lambda^*$ for the problem in (b).

    \item Is your problem in part (b) strongly dual? Prove that it is or isn't.

    \item Again for the specific model in part (b), give an argument for why the
      optimal dual variable $\lambda^*$ will always be strictly positive.

    \item Generate a random $5 \times 5$ $Q$ matrix in Python as $Q = A A^T$,
      where each $A_{ij} \sim N(0, 1)$. (Note: $Q$ will be positive-definite
      with probability 1.) Also, generate random $c_i \sim N(0,1)$. Solve the
      primal optimization problem defined by the quadratic constraint in part
      (b) using CVXPY. Print out the optimal point $x^*$ and optimal dual
      variable $\lambda^*$ and verify they match your analytic formulas from
      part (c) and that $\lambda^*$ is strictly positive.

    \item What is the value of $d^*$ (the optimal dual value) for your solution
      in part (f)? How does it compare to the optimal primal value $p^*$? Does
      this confirm your answer to part (d) about strong duality?

  \end{enumerate}

\item Exercise 5.14 in B\&V.

\item Exercise 5.37 in B\&V. You may want to refer to the following Wikipedia
  page for background for this and the next problem!

  \url{https://en.wikipedia.org/wiki/Discrete-time_Markov_chain}

\item A Markov chain is said to be reversible if there exists a probability vector $y$ such that
  $$p_{ij} y_j = p_{ji} y_i.$$

  \begin{enumerate}

  \item Show that if $P$ is reversible with respect to a probability vector $y$,
    then $y$ is a stationary distribution for $P$. That is,
    $$Py = y.$$

  \item Given any Markov chain transition matrix $P$ and any probability vector
    $y$, we can find the nearest reversible chain with respect to $y$ (in terms
    of Frobenius norm). Set this up as a convex optimization problem. ({\bf
      Hint:} As we've talked about in class, this will be easier if you
    formulate the objective function as a {\em squared} Frobenius norm.)

  \item Is this problem strongly dual? Prove that it is or isn't.

  \item Formulate the dual problem for the optimization problem you came up with
    in (b). Prove that the dual objective function $g(\lambda, \nu)$ is
    concave. (Don't just say its because $g(\lambda, \nu)$ is always concave!
    Show that the specific $g(\lambda, \nu)$ function you derived here is
    concave.)

  \item Solve your optimization problem in CVXPY for the specific case:
    $$P = \begin{pmatrix} 0 & 0 & 1\\ 1 & 0 & 0\\ 0 & 1 & 0\end{pmatrix}, \quad \text{and $y = (1/3, 1/3, 1/3)$}.$$
    Find additional solutions for probability vectors:
    $$y = (0.25, 0.25, 0.5), \quad \text{and} \quad y = (1, 0, 0).$$ For all
    three of the solutions, verify that the resulting optimal point satisfies
    the feasibility constraints. Also, verify that the corresponding $y$ is a
    stationary distribution (i.e., $P^* y = y$, where $P^*$ is the optimal
    transition matrix that your optimization found.)

  \item Print out the dual optimal points $\lambda^*$ and $\nu^*$ that are
    generated by CVXPY when you solve the primal problems with the same $P$ and
    three different $y$ values in (e).

  \item Now code up the dual problem as an optimization problem in CVXPY and
    solve it for the same three cases. Verify that the solutions for $\lambda^*$
    and $\nu*$ are the same as reported by CVXPY as dual optimal points in
    (f). After observing the results of this optimization, does this confirm
    your answer in part (c) about whether the problem is strongly dual? Why or
    why not?

  \item For each of the three $y$ values, rerun the optimization but perturb the
    $x \succeq 0$ constraint to instead be $x \succeq 0.01$. In each case, how
    does the optimal value change? Using the optimal dual variables, explain why
    each perturbed solution behaves the way it does. (Hint: one solution is
    particularly not like the others! Also, explain why it behaves so
    differently.)

  \end{enumerate}

\end{enumerate}

\end{document}
